\clearpage
\VerdierPhysicalPageAnchor{0125}%
\begingroup
\small
\setlength{\emergencystretch}{1.5em}%
\setlength{\abovedisplayskip}{4pt plus 2pt minus 2pt}%
\setlength{\belowdisplayskip}{4pt plus 2pt minus 2pt}%

\VerdierSemanticAnchor{ch2-1-3-5}
\noindent\textbf{Proposition 1.3.5.}\quad
\emph{Let $C$ be an abelian category. The functor}

\VerdierSemanticAnchor{ch2-1-3-5-h0-functor}
\[
  H^0\colon K(C)\longrightarrow C
\]
\noindent\emph{is a cohomological functor.}

\VerdierSemanticAnchor{ch2-1-3-5-first-exact-sequence}
Indeed, by (Chapter I, 3.3.1), it suffices to show that, for every object
$X\xrightarrow{f}Y$ of $\operatorname{tr}_1(C)$ (Chapter I, 3.1.1), the
sequence
\[
  H^0(X)\xrightarrow{H^0(f)}H^0(Y)
  \xrightarrow{H^0(p(f))}H^0(c(f))
  \xrightarrow{H^0(q(f))}H^0(X[1])
\]
\noindent is exact.

\VerdierSemanticAnchor{ch2-1-3-5-trii-reduction}
Using property TRII (1.1.1), one sees that it suffices to show that the
sequence

\VerdierSemanticAnchor{ch2-1-3-5-reduced-exact-sequence}
\[
  H^0(Y)\xrightarrow{H^0(p(f))}H^0(c(f))
  \xrightarrow{H^0(q(f))}H^0(X[1])
\]
\noindent is exact, which follows immediately from the definitions
(Chapter I, 3.1.2 and 3.2.2).

\VerdierSemanticAnchor{ch2-1-3-5-proof-complete}
The functor $H^0\colon K(C)\longrightarrow C$ is called the
\emph{canonical cohomological functor}.

\VerdierSemanticAnchor{ch2-1-3-5-shifted-cohomology}
In accordance with usage, put $H^n(X)=H^0(X[n])$.

\VerdierSemanticAnchor{ch2-1-3-6}
\medskip
\noindent\textbf{Proposition 1.3.6.}\quad
\emph{Let $C$ be an abelian category. The following conditions are
equivalent:}

\VerdierSemanticAnchor{ch2-1-3-6-equivalent-conditions}
\begin{enumerate}
\renewcommand{\labelenumi}{\roman{enumi})}
\item \emph{The category $K(C)$ is abelian.}
\item \emph{The category $K(C)$ is decomposable (1.2.8).}
\item \emph{The category $C$ is decomposable.}
\end{enumerate}

\VerdierSemanticAnchor{ch2-1-3-6-i-iff-ii}
The equivalence
$\mathrm{i)}\Longleftrightarrow\mathrm{ii)}$ follows from (1.2.9).

\VerdierSemanticAnchor{ch2-1-3-6-ii-implies-iii}
The implication $\mathrm{ii)}\Longrightarrow\mathrm{iii)}$ is clear because
$C$ is a full subcategory of $K(C)$.

\VerdierSemanticAnchor{ch2-1-3-6-hstar-complex}
For every complex $X$ in $C$, denote by $H^*(X)$ the complex
\[
  H^*(X)^l=H^l(X),\qquad l\in\mathbb Z,
\]
\[
  d_{H^*(X)}^l=0,\qquad l\in\mathbb Z.
\]

\VerdierSemanticAnchor{ch2-1-3-6-hstar-equivalence}
One checks that the functor
$H^*\colon K(C)\longrightarrow K(C)$ is an equivalence of categories when
$C$ is decomposable.

\VerdierSemanticAnchor{ch2-1-3-6-proof-complete}
It follows that
$\mathrm{iii)}\Longrightarrow\mathrm{ii)}$.
\endgroup
